% ============================================================================
% Fichier  : partie1/P01_C02_systeme_controle_interne.tex
% Titre    : Le Système de Contrôle Interne en Environnement Informatisé
% Auteur   : MINKA MI NGUIDJOÏ Thierry Emmanuel
% Version  : 0.1 -- novembre 2026
% ============================================================================
% Sources : COSO 2013, traduction IFACI (Normes) ; INF3226 C03 (resserré) ;
% décret n° 2013/287 (art. 28). INTOSAI GOV 9100 cité sans citation
% textuelle (texte absent du dossier Normes).
% Faits MINÉQUIR : profils de SIDÉPENSE (P00) ; matrice de C09 à garder
% cohérente : ENG (engagement), LIQ (liquidation), ORD (ordonnancement),
% FRS (fichier des fournisseurs), ADM (administration des comptes).
% ============================================================================

\chapter{Le Système de Contrôle Interne en Environnement Informatisé}
\label{chap:sci}
\chaptermark{Le système de contrôle interne}

\epigraph{%
    « L'organisation est-elle sous contrôle~? »%
}{Question à laquelle le Référentiel COSO veut permettre de répondre}

\niveauChapitre{Fondamental}{Chapitre~\ref{chap:audit-public}.}

\begin{boxobjectifs}
\begin{itemize}
    \item Définir le contrôle interne selon le COSO~2013 et en expliquer
          les quatre idées clés.
    \item Énoncer les trois catégories d'objectifs et les illustrer dans
          une administration financière.
    \item Présenter les cinq composantes et les dix-sept principes, et dire
          ce que chacun exige du système d'information.
    \item Classer un contrôle selon son moment (préventif, détectif,
          correctif) et son mode (manuel, automatisé, semi-automatisé).
    \item Distinguer les contrôles généraux informatiques des contrôles
          d'application, et expliquer pourquoi les premiers conditionnent
          les seconds.
    \item Traduire la séparation entre ordonnateur et comptable, et entre
          les étapes de la dépense, en profils d'une application.
    \item Énoncer les limites du contrôle interne.
\end{itemize}
\end{boxobjectifs}

\bigskip

% ============================================================================
\section*{Pourquoi ce chapitre}
% ============================================================================

Le prologue a posé une question~: qu'est-ce qui aurait dû empêcher le
double paiement~? La réponse porte un nom, le \emph{système de contrôle
interne}. Ce chapitre dit ce qu'il est~; le Chapitre~\ref{chap:eval-sci}
dit comment l'évaluer. Ensemble, ils forment le cœur du cours~: tous les
chapitres de la Partie~IV appliquent la même grille, domaine par domaine.

\separateur

% ============================================================================
\section{Qu'est-ce que le contrôle interne~?}
\label{sec:c02-definition}
% ============================================================================

\subsection{La définition du COSO}

Le cadre de référence le plus utilisé au monde est celui du \sigle{COSO}
(\textit{Committee of Sponsoring Organizations of the Treadway
Commission}), dans sa version de 2013, traduite en français par
l'IFACI~\cite{coso2013ic,coso2013icfr}.

\begin{boxdef}
\termecle{Contrôle interne}\index{contrôle interne}~: processus mis en
œuvre par l'organe de gouvernance, la direction et le personnel, destiné à
fournir une assurance raisonnable quant à la réalisation d'objectifs liés
aux opérations, à l'information (\textit{reporting}) et à la conformité.

\medskip
\termecle{Système de contrôle interne}\index{système de contrôle interne}
(\sigle{SCI})~: l'ensemble organisé des dispositifs de contrôle interne
d'une entité, considéré comme un tout~; c'est lui que l'auditeur évalue.
\end{boxdef}
{\small\itshape Définition du contrôle interne reformulée d'après le
COSO~2013.}

\medskip

La définition contient quatre idées, que l'auditeur public retrouve dans
son propre métier.
\begin{enumerate}
    \item \textbf{Un processus, pas un document.} Un manuel de procédures
          que personne n'applique n'est pas un contrôle interne.
    \item \textbf{L'affaire de tous.} Il est porté par l'ordonnateur, le
          chef de service, l'agent de saisie, l'informaticien, et pas
          seulement par le contrôleur.
    \item \textbf{Une assurance raisonnable.} Le coût d'un contrôle ne doit
          pas dépasser le risque qu'il traite~; aucun dispositif n'empêche
          tout.
    \item \textbf{Des objectifs.} On ne contrôle pas «~en général~», mais
          en vue d'objectifs précis~: sans objectif, pas de risque, et sans
          risque, pas de contrôle à justifier.
\end{enumerate}

\subsection{Trois catégories d'objectifs}

\begin{tableaucours}{L{2.6cm}L{4.6cm}L{5.0cm}}{%
    \tableauHeader{Catégorie} &
    \tableauHeader{Question} &
    \tableauHeader{Exemple au MINÉQUIR}
}
Opérations & Les activités sont-elles efficaces, les ressources
protégées~? & Chaque facture n'est payée qu'une fois, au bon
créancier \\
Information & L'information produite est-elle fiable~? & Le compte
administratif reflète les mandats réellement émis \\
Conformité & Les lois et règlements sont-ils respectés~? & Le circuit
engagement, liquidation, ordonnancement, paiement est respecté \\
\end{tableaucours}
\vspace{-0.8em}
\begin{center}{\small\itshape Tableau~2.1. Les objectifs du contrôle
interne}\end{center}

\subsection{Le contrôle interne dans le secteur public}

L'INTOSAI a adapté le cadre du COSO au secteur public dans ses lignes
directrices INTOSAI~GOV~9100. Elles ajoutent deux préoccupations propres à
l'administration~: la \textbf{protection des ressources publiques} contre
la perte, le gaspillage et l'usage abusif, et la conduite des opérations
dans le respect de l'\textbf{éthique} du service public. Le cours retient
le COSO~2013 comme cadre de travail, parce qu'il est le plus récent et le
seul dont le texte est dans le dossier des normes, et lit chaque principe
avec ces deux préoccupations.

Au Cameroun, l'évaluation des structures de contrôle interne est une
mission écrite du CONSUPE.

\begin{boxjuris}
\textbf{\loicam{Décret n°~2013/287}, article~28, alinéa~1}~\cite{decret2013287consupe}~:
les sections des divisions d'inspection sont chargées, chacune dans son
domaine de compétence, \textit{« de la préparation et de l'exécution des
opérations de vérification~; du suivi et de l'évaluation des activités
des structures de contrôle interne~; de la tenue de la documentation
technique et du fichier. »}

\medskip
\textbf{Portée pour l'audit.} Le vérificateur ne se contente pas de
relever des irrégularités~: il évalue les dispositifs qui auraient dû les
empêcher. L'article~27 place par ailleurs les systèmes d'information dans
le champ d'une section de la division chargée des administrations
publiques.
\end{boxjuris}

\separateur

% ============================================================================
\section{Cinq composantes, dix-sept principes}
\label{sec:c02-coso}
% ============================================================================

Le COSO organise le contrôle interne en \textbf{cinq composantes},
détaillées en \textbf{dix-sept principes}. Un dispositif n'est efficace
que si les cinq composantes sont mises en place, fonctionnent, et
fonctionnent conjointement~: une excellente activité de contrôle ne
compense pas une direction qui la contourne.

\begin{tableaucours}{L{3.2cm}cL{8.0cm}}{%
    \tableauHeader{Composante} &
    \tableauHeader{N°} &
    \tableauHeader{Principe (reformulé)}
}
Environnement de contrôle & 1 & S'engager en faveur de l'intégrité et de
l'éthique \\
 & 2 & Faire surveiller le contrôle interne par une instance indépendante
de la direction \\
 & 3 & Définir les structures, les pouvoirs et les responsabilités \\
 & 4 & Développer et retenir des personnes compétentes \\
 & 5 & Faire rendre compte chacun de ses responsabilités de contrôle \\
Évaluation des risques & 6 & Fixer des objectifs assez clairs pour en
identifier les risques \\
 & 7 & Identifier et analyser les risques \\
 & 8 & Évaluer le risque de fraude \\
 & 9 & Identifier et analyser les changements significatifs \\
Activités de contrôle & 10 & Choisir et mettre en place des activités de
contrôle qui réduisent les risques \\
 & 11 & Choisir et mettre en place des contrôles généraux informatiques \\
 & 12 & Déployer les contrôles par des politiques et des procédures \\
Information et communication & 13 & Produire et utiliser une information
pertinente et de qualité \\
 & 14 & Communiquer en interne sur le contrôle interne \\
 & 15 & Communiquer avec les tiers \\
Pilotage & 16 & Conduire des évaluations continues ou ponctuelles \\
 & 17 & Évaluer les déficiences et les communiquer à temps \\
\end{tableaucours}
\vspace{-0.8em}
\begin{center}{\small\itshape Tableau~2.2. Les composantes et les principes
du COSO~2013 (reformulation du cours)}\end{center}

\subsection{Les principes vus depuis le système d'information}

L'auditeur des SI lit chaque principe avec une question~: \emph{qu'est-ce
que ce principe exige du système d'information, et quelle pièce le
prouve~?} Le Tableau~2.3 en donne la réponse pour une administration
financière. C'est la base du questionnaire de contrôle interne du
Chapitre~\ref{chap:eval-sci}.

\begin{tableaucours}{cL{5.6cm}L{6.0cm}}{%
    \tableauHeader{N°} &
    \tableauHeader{Ce que le principe exige du SI} &
    \tableauHeader{Pièce que l'auditeur peut demander}
}
1 & Une charte d'usage des outils informatiques, connue des agents &
Charte, accusés de réception \\
2 & Une instance qui examine les risques informatiques & Comptes rendus du
comité de direction ou d'un comité dédié \\
3 & Des rôles écrits pour la cellule informatique et les propriétaires
d'application & Organigramme, décisions de nomination, délégations \\
4 & Des compétences informatiques entretenues & Plan de formation \\
5 & Des propriétaires d'application qui valident les droits & Revues des
habilitations signées \\
6 & Des objectifs écrits de disponibilité et de sécurité & Schéma
directeur, politique de sécurité \\
7 & Une analyse des risques informatiques & Cartographie des risques \\
8 & Une analyse des fraudes possibles par le système & Scénarios de
fraude (Chapitre~\ref{chap:donnees-fraude}) \\
9 & L'analyse des changements majeurs (nouvelle application, migration)
& Dossier de mise en service \\
10 & Des contrôles dans les applications & Paramétrage, contrôles
bloquants \\
11 & Des contrôles généraux informatiques & Accès, changements,
exploitation (section~\ref{sec:c02-cgi}) \\
12 & Des procédures informatiques écrites et appliquées & Procédures,
preuves d'exécution \\
13 & Des données fiables pour le pilotage & Rapprochements, contrôles de
qualité \\
14 & Une communication interne sur les règles & Notes de service \\
15 & Des échanges fiables avec les tiers (banques, trésor, prestataires)
& Conventions d'échange, contrats \\
16 & Une surveillance des contrôles & Tableaux de bord, revues
périodiques \\
17 & Un suivi des déficiences & Plan d'action, suivi des
recommandations \\
\end{tableaucours}
\vspace{-0.8em}
\begin{center}{\small\itshape Tableau~2.3. Les dix-sept principes vus
depuis le système d'information}\end{center}

\begin{boxinfo}
\textbf{Le principe~11 concerne directement l'auditeur des SI.} Le COSO
y demande de choisir et de mettre en place des contrôles généraux sur les
technologies~: les accès, les changements, l'exploitation. Les
Chapitres~\ref{chap:acces} et~\ref{chap:exploitation} les évaluent.
\end{boxinfo}

\separateur

% ============================================================================
\section{Classer les contrôles}
\label{sec:c02-types}
% ============================================================================

\subsection{Selon le moment}

\begin{boxdef}
\termecle{Contrôle préventif}\index{contrôle préventif}~: empêche
l'anomalie de se produire (SIDÉPENSE refuse un engagement supérieur au
crédit disponible).

\medskip
\termecle{Contrôle détectif}\index{contrôle détectif}~: repère l'anomalie
après coup (la revue mensuelle des mandats émis hors des heures de
service).

\medskip
\termecle{Contrôle correctif}\index{contrôle correctif}~: corrige
l'anomalie et ses effets (le recouvrement d'un paiement indu).
\end{boxdef}

Un dispositif solide combine les trois. Le préventif est préférable,
parce qu'il évite le dommage~; mais il peut être contourné, et seul le
détectif révèle alors le contournement. Au MINÉQUIR, tous les contrôles
préventifs du circuit ont manqué, et seuls deux contrôles détectifs ont
fonctionné (Chapitre~\ref{chap:prologue}).

\subsection{Selon le mode d'exécution}

\begin{tableaucours}{L{2.8cm}L{4.6cm}L{4.8cm}}{%
    \tableauHeader{Mode} &
    \tableauHeader{Forces} &
    \tableauHeader{Faiblesses}
}
Manuel (un visa) & Jugement humain, adaptable & Variable, oublié en fin
de gestion, preuve souvent absente \\
Automatisé (un blocage) & Constant, exhaustif & Ne vaut que si le
paramétrage est protégé~; une erreur se répète à chaque opération \\
Semi-automatisé (une liste d'anomalies examinée par un agent) & Combine
exhaustivité et jugement & Dépend de l'exactitude de la liste et de la
diligence de l'agent \\
\end{tableaucours}
\vspace{-0.8em}
\begin{center}{\small\itshape Tableau~2.4. Les modes d'exécution d'un
contrôle}\end{center}

\subsection{Selon la portée}

\begin{boxdef}
\termecle{Contrôles généraux informatiques}\index{contrôles généraux
informatiques} (\sigle{CGI}, \textit{IT general controls})~: contrôles
qui s'appliquent à l'environnement dans lequel les applications sont
exploitées et maintenues~: gestion des accès, gestion des changements,
exploitation (sauvegardes, traitements, incidents).

\medskip
\termecle{Contrôles d'application}\index{contrôles d'application}~:
contrôles intégrés à une application, qui agissent sur les opérations~:
contrôles de saisie, de traitement, d'interface et de sortie.
\end{boxdef}

\begin{boxalert}
\textbf{Violation typique.} Constater que SIDÉPENSE refuse tout mandat
supérieur au crédit disponible, et conclure que le contrôle budgétaire
est efficace. Or un compte d'administration peut modifier le crédit
directement dans la base, sans acte et sans trace.

\smallskip
\textbf{Règle~:} un contrôle d'application ne vaut que si les contrôles
généraux qui le protègent (accès, changements) sont efficaces.
L'auditeur les teste en premier.
\end{boxalert}

\separateur

% ============================================================================
\section{Les contrôles généraux informatiques}
\label{sec:c02-cgi}
% ============================================================================

\begin{tableaucours}{L{2.6cm}L{5.0cm}L{4.6cm}}{%
    \tableauHeader{Famille} &
    \tableauHeader{Objectif} &
    \tableauHeader{Contrôles typiques}
}
Accès & Seules les personnes autorisées accèdent aux applications et aux
données, avec les seuls droits nécessaires & Création et retrait des
comptes, revue des droits, comptes à privilèges \\
Changements & Aucune modification de programme ou de paramètre n'est mise
en service sans être demandée, testée et approuvée & Demande écrite,
tests, approbation, journal des mises en service \\
Exploitation & Les traitements s'exécutent comme prévu, les incidents sont
traités, les données sont sauvegardées & Planification, suivi des
incidents, sauvegardes testées \\
Acquisition et développement & Les nouvelles applications répondent aux
besoins et intègrent les contrôles & Cahier des charges, recette,
réception \\
\end{tableaucours}
\vspace{-0.8em}
\begin{center}{\small\itshape Tableau~2.5. Les quatre familles de
contrôles généraux informatiques}\end{center}

L'ordre compte. Si les accès ne sont pas maîtrisés, n'importe qui peut
modifier un paramètre~; si les changements ne le sont pas, un contrôle
automatique peut être désactivé sans que personne ne le sache~; dans les
deux cas, tester les contrôles d'application ne prouve plus rien. Le
cours suit donc cet ordre~: accès (Chapitre~\ref{chap:acces}),
changements et exploitation (Chapitre~\ref{chap:exploitation}), puis
contrôles d'application (Chapitre~\ref{chap:applications}).

\separateur

% ============================================================================
\section{Séparer les tâches dans une application}
\label{sec:c02-separation}
% ============================================================================

\subsection{Du principe au profil}

Le droit des finances publiques sépare l'ordonnateur, qui engage,
liquide et ordonnance, du comptable, qui paie après contrôle~; il sépare
aussi les étapes de la dépense~\cite{minfi2023regimefinancier}. Dans une
administration informatisée, ces séparations n'existent que si
l'application les impose.

\begin{boxdef}
\termecle{Séparation des tâches}\index{séparation des tâches}~: répartition
des étapes d'une opération entre des personnes différentes, de sorte
qu'aucune ne puisse à elle seule réaliser et dissimuler une irrégularité.

\medskip
\termecle{Profil}\index{profil}~: ensemble de droits attribué dans une
application (créer un engagement, liquider, ordonnancer, modifier un
RIB). Un utilisateur peut avoir plusieurs profils~: c'est leur
combinaison qui crée, ou non, un conflit.
\end{boxdef}

\subsection{La matrice des profils de SIDÉPENSE}

SIDÉPENSE connaît cinq profils~: ENG (engagement), LIQ (liquidation), ORD
(ordonnancement), FRS (gestion du fichier des fournisseurs, dont les RIB)
et ADM (administration des comptes). Le Tableau~2.6 indique les
combinaisons qui ne doivent jamais être réunies sur un même compte.

\begin{tableaucours}{L{2.2cm}L{4.4cm}L{5.6cm}}{%
    \tableauHeader{Paire} &
    \tableauHeader{Pourquoi elle est interdite} &
    \tableauHeader{Ce qu'elle permet}
}
LIQ + ORD & L'ordonnancement contrôle la liquidation & Arrêter un montant et
émettre seul le mandat \\
FRS + ORD & Le RIB doit être fixé par un autre que celui qui émet le
mandat & Changer le bénéficiaire puis le payer \\
FRS + LIQ & Même raison, en amont & Créer un service fait au profit d'un
compte choisi \\
ADM + tout profil métier & L'administrateur peut s'attribuer n'importe
quel droit & Contourner toutes les séparations \\
\end{tableaucours}
\vspace{-0.8em}
\begin{center}{\small\itshape Tableau~2.6. Les paires de profils
incompatibles dans SIDÉPENSE}\end{center}

\begin{boxscenario}{Dossier MINÉQUIR. Le compte mbarga.j}
\textbf{Constatation.} À la date du 28~décembre~2025, le compte
\texttt{mbarga.j} détient les profils LIQ, ORD et FRS (matrice des
profils extraite le 23~mars~2026).

\medskip
\textbf{Analyse} (confiance élevée). Le compte réunit les trois paires
interdites du Tableau~2.6 qui ne demandent pas l'administration. Un seul
compte suffisait pour modifier le RIB, liquider et ordonnancer~: le
contrôle interne du circuit reposait entièrement sur la confiance en une
personne, et cette personne n'était plus en poste.
\end{boxscenario}

\subsection{Quand on ne peut pas séparer}

Dans un petit service, trois agents ne peuvent pas tenir cinq profils
sans cumul. La pratique admet alors des \termecle{contrôles
compensatoires}\index{contrôle compensatoire}~: une revue, par une
personne indépendante, de toutes les opérations réalisées par un compte
qui cumule des profils. Le contrôle compensatoire doit être écrit,
exécuté et prouvé~; sinon, le cumul est une déficience.

\separateur

% ============================================================================
\section{Les limites du contrôle interne}
\label{sec:c02-limites}
% ============================================================================

Le COSO rappelle que le contrôle interne, même bien conçu, a des
limites~\cite{coso2013icfr}. L'auditeur doit les connaître, pour ne pas
reprocher à un dispositif ce qu'aucun dispositif ne peut empêcher, et pour
savoir où chercher.

\begin{itemize}
    \item \textbf{Le jugement humain} peut être erroné~; les décisions
          sont prises sous la contrainte du temps.
    \item \textbf{Les défaillances} existent~: fatigue, erreur, oubli,
          surtout en fin de gestion.
    \item \textbf{La collusion} de deux personnes ou plus défait la
          séparation des tâches.
    \item \textbf{Le contournement par la direction}~: un responsable peut
          passer outre un contrôle qu'il a lui-même mis en place.
    \item \textbf{Les événements extérieurs} échappent à l'entité.
\end{itemize}

\begin{boxmetier}{Contrôleur financier}
Le contrôle a priori que vous exercez est une activité de contrôle au sens
du COSO (principe~10). S'il passe par SIDÉPENSE, demandez à voir le profil
qui porte votre visa, et la liste des personnes qui le détiennent~: votre
visa ne vaut que si personne d'autre ne peut l'apposer.
\end{boxmetier}

\begin{boxmetier}{Auditeur interne}
Le COSO fait du pilotage (principes~16 et~17) une composante du contrôle
interne~: votre propre travail en fait partie. Gardez une liste des
déficiences que vous avez signalées et de leur suivi~; le vérificateur la
demandera.
\end{boxmetier}

\separateur

\begin{boxresume}
\begin{itemize}
    \item \textbf{Un processus orienté vers des objectifs.} Le contrôle
          interne donne une assurance raisonnable sur les opérations,
          l'information et la conformité.
    \item \textbf{Cinq composantes, dix-sept principes.} Environnement,
          risques, activités de contrôle, information et communication,
          pilotage~; chacune doit être en place, fonctionner, et
          fonctionner avec les autres.
    \item \textbf{Chaque principe a une traduction informatique}, et une
          pièce qui la prouve.
    \item \textbf{Préventif, détectif, correctif~; manuel, automatisé,
          semi-automatisé.} Un dispositif solide combine les modes.
    \item \textbf{Les contrôles généraux d'abord.} Sans maîtrise des accès
          et des changements, aucun contrôle d'application n'est fiable.
    \item \textbf{La séparation se lit dans les profils.} Elle n'existe que
          si l'application l'impose~; sinon, il faut un contrôle
          compensatoire prouvé.
    \item \textbf{Une mission écrite du CONSUPE.} Le décret n°~2013/287
          (art.~28) le charge d'évaluer les structures de contrôle
          interne.
\end{itemize}
\end{boxresume}

\begin{boxexo}[title={\ding{45}\quad Exercice~2.1. Classer des contrôles \niveaubase}]
Classez chaque contrôle selon son moment et son mode.
\begin{enumerate}[(a)]
    \item SIDÉPENSE refuse un engagement supérieur au crédit disponible.
    \item Le chef de la liquidation signe chaque soir la liste des mandats
          du jour, éditée par l'application.
    \item Le poste comptable rejette un mandat dont la facture scannée est
          illisible.
    \item La DAF émet un ordre de reversement après un paiement indu.
\end{enumerate}
\end{boxexo}

\begin{boxexo}[title={\ding{45}\quad Exercice~2.2. Composantes et principes \niveaubase}]
Rattachez chacun des manques suivants, relevés au MINÉQUIR, à une
composante et à un principe du COSO~: (a) aucune cartographie des risques
de SIDÉPENSE~; (b) aucune procédure de désactivation des comptes lors des
mutations~; (c) les anomalies de fin de gestion ne remontent jamais au
directeur~; (d) la cellule informatique n'a aucune fiche de poste.
\end{boxexo}

\begin{boxexo}[title={\ding{45}\quad Exercice~2.3. Une matrice des profils \niveaumoyen}]
La DAF compte six utilisateurs de SIDÉPENSE~: ABENA (ENG), BELLO (LIQ,
ORD), CHOUPO (LIQ), DJOMO (ORD, FRS), ESSOMBA (FRS), FOUDA (ADM).
\begin{enumerate}[(a)]
    \item Identifiez les comptes en conflit au regard du Tableau~2.6.
    \item Proposez une nouvelle répartition sans conflit, sans créer de
          compte.
    \item Si une répartition sans conflit est impossible, décrivez le
          contrôle compensatoire et la preuve qu'il doit laisser.
\end{enumerate}
\end{boxexo}

\begin{boxexo}[title={\ding{45}\quad Exercice~2.4. Les contrôles généraux d'abord \niveauexpert}]
Le chef de la cellule informatique vous montre que SIDÉPENSE bloque
désormais les doublons (même fournisseur, même montant, même facture).
Dites quels contrôles généraux vous devez tester avant de conclure que ce
blocage est efficace pour toute l'année 2026, quelle pièce vous demandez
pour chacun, et ce que vous concluez si l'une de ces pièces manque.
\end{boxexo}

\begin{boxinfo}
\textbf{Transition.} Le contrôle interne est maintenant défini, composante
par composante. Reste la question de l'auditeur~: comment savoir s'il est
efficace~? Le Chapitre~\ref{chap:eval-sci} en donne la démarche, du
questionnaire à la conclusion.
\end{boxinfo}
